% --- Archon physics formalization source begin ---
% archon:physics
% archon:covers PhyXMiniProblems/problem_phyx_mini_0367.lean
% archon:source-report reports/phyx_mini/problem_phyx_mini_0367.source.json

\chapter{Physics problem phyx\_mini\_0367}
\label{ch:PhyXMiniProblems_problem_phyx_mini_0367}

\paragraph{Problem source.}
\#\# Physical scenario

8.0\textbackslash{},\textbackslash{}text\{g\} of helium gas follows the process \textbackslash{}( 1 \textbackslash{}rightarrow 2 \textbackslash{}rightarrow 3 \textbackslash{}) shown in figure.

\#\# Question

Find the value of \textbackslash{}( T\_3 \textbackslash{}).

\#\# Answer choices

- A: A: 657\textasciicircum{}\textbackslash{}circ C

- B: B: 556\textasciicircum{}\textbackslash{}circ C

- C: C: 636\textasciicircum{}\textbackslash{}circ C

- D: D: 784\textasciicircum{}\textbackslash{}circ C

\#\# Figure caption (auxiliary; use the image as primary evidence)

The image is a graph depicting a thermodynamic process on a pressure-volume (p-V) diagram. Here's the fine-grained content:

- **Axes**: 
  - The vertical axis is labeled \textbackslash{}( p \textbackslash{}) (pressure) and shows a value of \textbackslash{}( 2 \textbackslash{}, \textbackslash{}text\{atm\} \textbackslash{}).
  - The horizontal axis is labeled \textbackslash{}( V \textbackslash{}) (volume).

- **Points and Path**:
  - Three points are marked: 1, 2, and 3.
  - Point 1 to 2: A vertical line from point 1 to point 2. 
    - Point 1 is labeled as \textbackslash{}( 37\textasciicircum{}\textbackslash{}circ\textbackslash{}text\{C\} \textbackslash{}).
    - Point 2 is labeled as \textbackslash{}( 657\textasciicircum{}\textbackslash{}circ\textbackslash{}text\{C\} \textbackslash{}) and is above \textbackslash{}( p\_2 \textbackslash{}).
  - Point 2 to 3: A curved line from point 2 to point 3, labeled "Isothermal."

- **Labels**:
  - \textbackslash{}( p\_2 \textbackslash{}) is marked on the vertical axis, corresponding to the pressure at point 2.
  - Volumes \textbackslash{}( V\_1 \textbackslash{}) and \textbackslash{}( V\_3 \textbackslash{}) are marked on the horizontal axis.
  
The diagram represents a thermodynamic process moving from point 1 to 2 isobarically (constant pressure), and then from point 2 to 3 isothermally (constant temperature).

\#\# Dataset metadata

Ideal Gases and Kinetic Theory; Physical Model Grounding Reasoning, Multi-Formula Reasoning

\paragraph{Recorded answer/context.}
A: A: 657\textasciicircum{}\textbackslash{}circ C

\paragraph{Figure/image path.}
/root/proposal\_for\_physic/science-mango/phyx\_mini\_run/phyx\_data/test\_image/367.png

\paragraph{Formalization target.}
create a compiling Lean file with sorry bodies at `PhyXMiniProblems/problem\_phyx\_mini\_0367.lean`. The Lean declarations must preserve the physical quantities, dimensions or dimensional roles, figure labels, governing-law hypotheses, and final relation expressed by this problem.
Use Mathlib/Physlib names found through LeanExplore where available. If a physics API is missing, introduce faithful local abstractions rather than scalar placeholder aliases.

\begin{theorem}[Physics formalization target]
\label{thm:physics:phyx_mini_0367:target}
\lean{PhyXMiniProblems.ProblemPhyXMini0367.temperatureAtStateThree_eq_657_celsius}
\uses{def:physics:phyx-mini-0367:phyxminiproblems-problemphyxmini0367-temperatureincelsius, def:physics:phyx-mini-0367:phyxminiproblems-problemphyxmini0367-heliumpvprocesssetup, def:physics:phyx-mini-0367:phyxminiproblems-problemphyxmini0367-hasphysicalparameters, def:physics:phyx-mini-0367:phyxminiproblems-problemphyxmini0367-satisfiesprocesssemantics, def:physics:phyx-mini-0367:phyxminiproblems-problemphyxmini0367-hasscenarioandfiguredata}
The curved leg \(2\to3\) is isothermal, so state 3 has the same
\(657^\circ\mathrm C\) temperature printed at state 2.
\end{theorem}
\begin{proof}
The figure data identify leg \(2\to3\) as the isothermal leg.  Applying the
process law to that leg gives equality of the two dimensionful absolute
temperatures.  Applying the Celsius readout to this equality and substituting
the printed state-2 value yields
\(T_3=657^\circ\mathrm C\).
\end{proof}

% --- Archon declaration topology begin ---
\section{Declaration topology}
\begin{definition}[Mass Quantity]
  \label{def:physics:phyx-mini-0367:phyxminiproblems-problemphyxmini0367-massquantity}
  \lean{PhyXMiniProblems.ProblemPhyXMini0367.MassQuantity}
  A nonnegative physical mass, of dimension mass
\end{definition}
\begin{proof}
  This is the declared model definition of Mass Quantity.
\end{proof}
\begin{definition}[Volume Quantity]
  \label{def:physics:phyx-mini-0367:phyxminiproblems-problemphyxmini0367-volumequantity}
  \lean{PhyXMiniProblems.ProblemPhyXMini0367.VolumeQuantity}
  A nonnegative physical volume, of dimension length³
\end{definition}
\begin{proof}
  This is the declared model definition of Volume Quantity.
\end{proof}
\begin{definition}[mass In Grams]
  \label{def:physics:phyx-mini-0367:phyxminiproblems-problemphyxmini0367-massingrams}
  \lean{PhyXMiniProblems.ProblemPhyXMini0367.massInGrams}
  \uses{def:physics:phyx-mini-0367:phyxminiproblems-problemphyxmini0367-massquantity}
  Read a physical mass in grams
\end{definition}
\begin{proof}
  This is the declared model definition of mass In Grams.
\end{proof}
\begin{definition}[volume In Cubic Meters]
  \label{def:physics:phyx-mini-0367:phyxminiproblems-problemphyxmini0367-volumeincubicmeters}
  \lean{PhyXMiniProblems.ProblemPhyXMini0367.volumeInCubicMeters}
  \uses{def:physics:phyx-mini-0367:phyxminiproblems-problemphyxmini0367-volumequantity}
  Read a physical volume in SI cubic metres
\end{definition}
\begin{proof}
  This is the declared model definition of volume In Cubic Meters.
\end{proof}
\begin{definition}[pressure In Pascals]
  \label{def:physics:phyx-mini-0367:phyxminiproblems-problemphyxmini0367-pressureinpascals}
  \lean{PhyXMiniProblems.ProblemPhyXMini0367.pressureInPascals}
  Read a physical pressure in SI pascals
\end{definition}
\begin{proof}
  This is the declared model definition of pressure In Pascals.
\end{proof}
\begin{definition}[temperature In Kelvin]
  \label{def:physics:phyx-mini-0367:phyxminiproblems-problemphyxmini0367-temperatureinkelvin}
  \lean{PhyXMiniProblems.ProblemPhyXMini0367.temperatureInKelvin}
  Physlib's Temperature has an arbitrary absolute scale. For this problem that scale is calibrated in kelvin, so Temperature.toReal is the kelvin readout
\end{definition}
\begin{proof}
  This is the declared model definition of temperature In Kelvin.
\end{proof}
\begin{definition}[temperature In Celsius]
  \label{def:physics:phyx-mini-0367:phyxminiproblems-problemphyxmini0367-temperatureincelsius}
  \lean{PhyXMiniProblems.ProblemPhyXMini0367.temperatureInCelsius}
  \uses{def:physics:phyx-mini-0367:phyxminiproblems-problemphyxmini0367-temperatureinkelvin}
  Celsius readout obtained from the calibrated absolute temperature
\end{definition}
\begin{proof}
  This is the declared model definition of temperature In Celsius.
\end{proof}
\begin{definition}[Process State]
  \label{def:physics:phyx-mini-0367:phyxminiproblems-problemphyxmini0367-processstate}
  \lean{PhyXMiniProblems.ProblemPhyXMini0367.ProcessState}
  The three equilibrium states labelled 1, 2, and 3 in the image
\end{definition}
\begin{proof}
  This is the declared model definition of Process State.
\end{proof}
\begin{definition}[Process Leg]
  \label{def:physics:phyx-mini-0367:phyxminiproblems-problemphyxmini0367-processleg}
  \lean{PhyXMiniProblems.ProblemPhyXMini0367.ProcessLeg}
  The two directed process legs shown by arrows in the image
\end{definition}
\begin{proof}
  This is the declared model definition of Process Leg.
\end{proof}
\begin{definition}[initial State]
  \label{def:physics:phyx-mini-0367:phyxminiproblems-problemphyxmini0367-processleg-initialstate}
  \lean{PhyXMiniProblems.ProblemPhyXMini0367.ProcessLeg.initialState}
  \uses{def:physics:phyx-mini-0367:phyxminiproblems-problemphyxmini0367-processstate, def:physics:phyx-mini-0367:phyxminiproblems-problemphyxmini0367-processleg}
  Initial state of a directed process leg
\end{definition}
\begin{proof}
  This is the declared model definition of initial State.
\end{proof}
\begin{definition}[final State]
  \label{def:physics:phyx-mini-0367:phyxminiproblems-problemphyxmini0367-processleg-finalstate}
  \lean{PhyXMiniProblems.ProblemPhyXMini0367.ProcessLeg.finalState}
  \uses{def:physics:phyx-mini-0367:phyxminiproblems-problemphyxmini0367-processstate, def:physics:phyx-mini-0367:phyxminiproblems-problemphyxmini0367-processleg}
  Final state of a directed process leg
\end{definition}
\begin{proof}
  This is the declared model definition of final State.
\end{proof}
\begin{definition}[Process Kind]
  \label{def:physics:phyx-mini-0367:phyxminiproblems-problemphyxmini0367-processkind}
  \lean{PhyXMiniProblems.ProblemPhyXMini0367.ProcessKind}
  Thermodynamic character of a process leg
\end{definition}
\begin{proof}
  This is the declared model definition of Process Kind.
\end{proof}
\begin{definition}[Working Substance]
  \label{def:physics:phyx-mini-0367:phyxminiproblems-problemphyxmini0367-workingsubstance}
  \lean{PhyXMiniProblems.ProblemPhyXMini0367.WorkingSubstance}
  Working substance named in the scenario
\end{definition}
\begin{proof}
  This is the declared model definition of Working Substance.
\end{proof}
\begin{definition}[Figure Axis]
  \label{def:physics:phyx-mini-0367:phyxminiproblems-problemphyxmini0367-figureaxis}
  \lean{PhyXMiniProblems.ProblemPhyXMini0367.FigureAxis}
  Horizontal or vertical axis of the supplied diagram
\end{definition}
\begin{proof}
  This is the declared model definition of Figure Axis.
\end{proof}
\begin{definition}[Figure Axis Quantity]
  \label{def:physics:phyx-mini-0367:phyxminiproblems-problemphyxmini0367-figureaxisquantity}
  \lean{PhyXMiniProblems.ProblemPhyXMini0367.FigureAxisQuantity}
  Physical quantity assigned to a figure axis
\end{definition}
\begin{proof}
  This is the declared model definition of Figure Axis Quantity.
\end{proof}
\begin{definition}[PVSegment Shape]
  \label{def:physics:phyx-mini-0367:phyxminiproblems-problemphyxmini0367-pvsegmentshape}
  \lean{PhyXMiniProblems.ProblemPhyXMini0367.PVSegmentShape}
  Qualitative geometric shape of a process segment in the raster
\end{definition}
\begin{proof}
  This is the declared model definition of PVSegment Shape.
\end{proof}
\begin{definition}[Thermodynamic State Data]
  \label{def:physics:phyx-mini-0367:phyxminiproblems-problemphyxmini0367-thermodynamicstatedata}
  \lean{PhyXMiniProblems.ProblemPhyXMini0367.ThermodynamicStateData}
  \uses{def:physics:phyx-mini-0367:phyxminiproblems-problemphyxmini0367-volumequantity}
  Pressure, volume, and absolute temperature at one equilibrium state
\end{definition}
\begin{proof}
  This is the declared model definition of Thermodynamic State Data.
\end{proof}
\begin{definition}[Pressure Volume Figure]
  \label{def:physics:phyx-mini-0367:phyxminiproblems-problemphyxmini0367-pressurevolumefigure}
  \lean{PhyXMiniProblems.ProblemPhyXMini0367.PressureVolumeFigure}
  \uses{def:physics:phyx-mini-0367:phyxminiproblems-problemphyxmini0367-processstate, def:physics:phyx-mini-0367:phyxminiproblems-problemphyxmini0367-processleg, def:physics:phyx-mini-0367:phyxminiproblems-problemphyxmini0367-figureaxis, def:physics:phyx-mini-0367:phyxminiproblems-problemphyxmini0367-figureaxisquantity, def:physics:phyx-mini-0367:phyxminiproblems-problemphyxmini0367-pvsegmentshape}
  Qualitative evidence retained from the supplied raster. Coordinates record only relative placement in the drawing; physical state quantities are stored separately in HeliumPVProcessSetup
\end{definition}
\begin{proof}
  This is the declared model definition of Pressure Volume Figure.
\end{proof}
\begin{definition}[Helium PVProcess Setup]
  \label{def:physics:phyx-mini-0367:phyxminiproblems-problemphyxmini0367-heliumpvprocesssetup}
  \lean{PhyXMiniProblems.ProblemPhyXMini0367.HeliumPVProcessSetup}
  \uses{def:physics:phyx-mini-0367:phyxminiproblems-problemphyxmini0367-massquantity, def:physics:phyx-mini-0367:phyxminiproblems-problemphyxmini0367-volumequantity, def:physics:phyx-mini-0367:phyxminiproblems-problemphyxmini0367-processstate, def:physics:phyx-mini-0367:phyxminiproblems-problemphyxmini0367-processleg, def:physics:phyx-mini-0367:phyxminiproblems-problemphyxmini0367-processkind, def:physics:phyx-mini-0367:phyxminiproblems-problemphyxmini0367-workingsubstance, def:physics:phyx-mini-0367:phyxminiproblems-problemphyxmini0367-thermodynamicstatedata, def:physics:phyx-mini-0367:phyxminiproblems-problemphyxmini0367-pressurevolumefigure}
  Independent physical quantities in the scenario, including the labels V₁, V₃, and p₂. No field fixes the requested temperature at state 3
\end{definition}
\begin{proof}
  This is the declared model definition of Helium PVProcess Setup.
\end{proof}
\begin{definition}[Has Physical Parameters]
  \label{def:physics:phyx-mini-0367:phyxminiproblems-problemphyxmini0367-hasphysicalparameters}
  \lean{PhyXMiniProblems.ProblemPhyXMini0367.HasPhysicalParameters}
  \uses{def:physics:phyx-mini-0367:phyxminiproblems-problemphyxmini0367-massingrams, def:physics:phyx-mini-0367:phyxminiproblems-problemphyxmini0367-volumeincubicmeters, def:physics:phyx-mini-0367:phyxminiproblems-problemphyxmini0367-pressureinpascals, def:physics:phyx-mini-0367:phyxminiproblems-problemphyxmini0367-temperatureinkelvin, def:physics:phyx-mini-0367:phyxminiproblems-problemphyxmini0367-heliumpvprocesssetup}
  Positivity conditions selecting physical pressures, volumes, and temperatures
\end{definition}
\begin{proof}
  This is the declared model definition of Has Physical Parameters.
\end{proof}
\begin{definition}[Satisfies Process Semantics]
  \label{def:physics:phyx-mini-0367:phyxminiproblems-problemphyxmini0367-satisfiesprocesssemantics}
  \lean{PhyXMiniProblems.ProblemPhyXMini0367.SatisfiesProcessSemantics}
  \uses{def:physics:phyx-mini-0367:phyxminiproblems-problemphyxmini0367-heliumpvprocesssetup}
  Governing semantics of the two named process kinds. These are general laws: an isochoric leg preserves physical volume, and an isothermal leg preserves absolute temperature. They contain no numerical value for T₃
\end{definition}
\begin{proof}
  This is the declared model definition of Satisfies Process Semantics.
\end{proof}
\begin{definition}[Has Scenario And Figure Data]
  \label{def:physics:phyx-mini-0367:phyxminiproblems-problemphyxmini0367-hasscenarioandfiguredata}
  \lean{PhyXMiniProblems.ProblemPhyXMini0367.HasScenarioAndFigureData}
  \uses{def:physics:phyx-mini-0367:phyxminiproblems-problemphyxmini0367-massingrams, def:physics:phyx-mini-0367:phyxminiproblems-problemphyxmini0367-temperatureincelsius, def:physics:phyx-mini-0367:phyxminiproblems-problemphyxmini0367-heliumpvprocesssetup}
  Scenario data and direct readouts from the primary image. In particular, states 1 and 3 lie at the 2 atm height, states 1 and 2 lie over V₁, state 3 lies over V₃, and the image labels the temperature of state 2 as 657 °C. The requested state-3 temperature is deliberately absent
\end{definition}
\begin{proof}
  This is the declared model definition of Has Scenario And Figure Data.
\end{proof}
% --- Archon declaration topology end ---
% --- Archon physics formalization source end ---
